\documentclass[12pt,a4paper]{article}
\usepackage[utf8]{inputenc}
\usepackage[T1,T2A]{fontenc}
\usepackage[russian]{babel}
\usepackage{amsmath,amssymb}
\usepackage[margin=2.5cm]{geometry}

\begin{document}

\section*{Формулировка задачи поиска схемы загрузки активной зоны в форме математического программирования}

\subsection*{Управляемые переменные}

Пусть активная зона реактора содержит $N$ позиций в секторе симметрии 1/3. Введём следующие управляемые переменные:

\textbf{Переменная 1 --- выбор шаблона:}

\begin{equation}
T \in \mathcal{T}
\end{equation}

где $\mathcal{T}$ --- заданное конечное множество допустимых шаблонов.

\textbf{Определение шаблона.} Шаблон $T$ представляет собой бинарную матрицу

\begin{equation}
T = (\tau_{bi}) \in \{0,1\}^{3 \times N}, \quad b \in \{0,1,2\},\; i = 1,\ldots,N
\end{equation}

где $\tau_{bi} = 1$, если позиция $i$ в шаблоне отнесена к типу $b$ (0 --- свежая, 1 --- однажды выгоревшая, 2 --- дважды выгоревшая), и $\tau_{bi} = 0$ в противном случае.

Поскольку каждая позиция должна иметь ровно один тип, матрица $T$ обязана удовлетворять условию:

\begin{equation}
\sum_{b=0}^{2} \tau_{bi} = 1, \quad \forall\, i = 1,\ldots,N
\end{equation}

Шаблон $T$ однозначно определяет разбиение позиций сектора на три непересекающихся подмножества:

\begin{equation}
\mathcal{P}_b(T) = \{\, i \,:\, \tau_{bi} = 1 \,\}, \quad b \in \{0,1,2\}
\end{equation}

\begin{equation}
\mathcal{P}_0(T) \cup \mathcal{P}_1(T) \cup \mathcal{P}_2(T) = \mathcal{P}, \quad \mathcal{P}_k(T) \cap \mathcal{P}_l(T) = \varnothing \;\; (k \neq l)
\end{equation}

\begin{equation}
|\mathcal{P}_k(T)| = N_k, \quad N_0 + N_1 + N_2 = N
\end{equation}

--- позиции для свежих, однажды и дважды выгоревших ТВС соответственно.

Множество допустимых шаблонов $\mathcal{T}$ --- это заранее заданный (например, отобранный экспертным путём) список таких матриц, удовлетворяющих указанному условию и дополнительным физическим требованиям (соотношение $N_0, N_1, N_2$).

\textbf{Переменная 2 --- назначение видов свежих ТВС на позиции:}

\begin{equation}
\sigma_0 : \mathcal{P}_0 \to \mathcal{F}
\end{equation}

где $\mathcal{F} = \{f_1, \ldots, f_M\}$ --- множество доступных видов свежих ТВС (виды различаются обогащением по $^{235}\text{U}$, количеством твэг с выгорающим поглотителем и его концентрацией). Функция $\sigma_0$ сопоставляет каждой позиции $i \in \mathcal{P}_0$ вид свежей ТВС $\sigma_0(i) \in \mathcal{F}$. Одному виду может соответствовать несколько позиций.

\textbf{Переменная 3 --- перестановка однажды выгоревших ТВС:}

\begin{equation}
\sigma_1 : \mathcal{P}_1 \to \mathcal{U}_1
\end{equation}

где $\mathcal{U}_1 = \{u_1^{(1)}, \ldots, u_{N_0}^{(1)}\}$ --- набор конкретных однажды выгоревших ТВС из предыдущего цикла ($|\mathcal{U}_1| = N_0$). Функция $\sigma_1$ является инъекцией: каждая ТВС из $\mathcal{U}_1$ устанавливается не более чем в одну позицию.

\textbf{Переменная 4 --- перестановка дважды выгоревших ТВС:}

\begin{equation}
\sigma_2 : \mathcal{P}_2 \to \mathcal{U}_2
\end{equation}

где $\mathcal{U}_2 = \{u_1^{(2)}, \ldots, u_{N_0}^{(2)}\}$ --- набор конкретных дважды выгоревших ТВС ($|\mathcal{U}_2| = N_0$). Функция $\sigma_2$ также является инъекцией.

Совокупное решение задаётся четвёркой:

\begin{equation}
\mathbf{s} = (T,\; \sigma_0,\; \sigma_1,\; \sigma_2)
\end{equation}

\subsection*{Целевая функция}

Требуется найти $\mathbf{s}$, доставляющее максимум функции:

\begin{equation}
\max_{\mathbf{s}}\; F(\mathbf{s}) = -\mathrm{LCOE}(\mathbf{s}) - \sum_{c \in \mathcal{C}} \gamma_c \cdot \Phi\left(v_c(\mathbf{s}),\; \bar{v}_c\right) + 1 \times \delta_{v_c(\mathbf{s}) \leq \bar{v}_c}
\end{equation}

где:
\begin{itemize}
\item $\mathrm{LCOE}(\mathbf{s})$ --- нормированная стоимость выработки электроэнергии за топливный цикл;
\item $\mathcal{C}$ --- множество ограничивающих физических параметров ($|\mathcal{C}| \sim 5$--$10$);
\item $\gamma_c > 0$ --- весовой коэффициент для параметра $c$;
\item $v_c(\mathbf{s})$ --- значение $c$-го физического параметра при схеме $\mathbf{s}$, вычисляемое расчётным кодом;
\item $\bar{v}_c$ --- допустимый предел по параметру $c$;
\item $\Phi(v_c(\mathbf{s}), \bar{v}_c)$ --- штрафная функция нарушения ограничения:
\end{itemize}

\begin{equation}
\Phi\left(v_c(\mathbf{s}),\; \bar{v}_c\right) = \delta_{v_c(\mathbf{s}) > \bar{v}_c} \cdot \left(\frac{v_c(\mathbf{s}) - \bar{v}_c}{\bar{v}_c}\right)^2
\end{equation}

где $\delta_{v_c > \bar v_c}$ --- индикаторная функция (равна 1, если условие выполняется, иначе 0).

\subsection*{Система ограничений}

\textbf{Г1 --- Выбор шаблона из допустимого множества:}

\begin{equation}
T \in \mathcal{T}
\end{equation}

\textbf{Г2 --- Каждая позиция для свежих ТВС получает ровно один вид ТВС:}

\begin{equation}
\sigma_0(i) \in \mathcal{F}, \quad \forall\, i \in \mathcal{P}_0(T)
\end{equation}

\textbf{Г3 --- Каждая позиция для однажды выгоревших ТВС получает ровно одну конкретную ТВС:}

\begin{equation}
\sigma_1(i) \in \mathcal{U}_1, \quad \forall\, i \in \mathcal{P}_1(T)
\end{equation}

\textbf{Г4 --- Каждая позиция для дважды выгоревших ТВС получает ровно одну конкретную ТВС:}

\begin{equation}
\sigma_2(i) \in \mathcal{U}_2, \quad \forall\, i \in \mathcal{P}_2(T)
\end{equation}

\textbf{Г5 --- Инъективность назначения однажды выгоревших ТВС} (одна ТВС не может стоять в двух позициях одновременно):

\begin{equation}
\sigma_1(i) \neq \sigma_1(i'), \quad \forall\, i \neq i',\quad i, i' \in \mathcal{P}_1(T)
\end{equation}

\textbf{Г6 --- Инъективность назначения дважды выгоревших ТВС:}

\begin{equation}
\sigma_2(i) \neq \sigma_2(i'), \quad \forall\, i \neq i',\quad i, i' \in \mathcal{P}_2(T)
\end{equation}

\textbf{Г7 --- Область определения решения:}

\begin{equation}
\mathbf{s} = (T,\; \sigma_0,\; \sigma_1,\; \sigma_2), \quad T \in \mathcal{T},\quad \sigma_0 \in \mathcal{F}^{\mathcal{P}_0},\quad \sigma_1 \in \mathrm{Inj}(\mathcal{P}_1, \mathcal{U}_1),\quad \sigma_2 \in \mathrm{Inj}(\mathcal{P}_2, \mathcal{U}_2)
\end{equation}

где $\mathrm{Inj}(\mathcal{P}_k, \mathcal{U}_k)$ --- множество всех инъекций из $\mathcal{P}_k$ в $\mathcal{U}_k$.

\subsection*{Итоговая запись задачи}

\begin{equation}
\max_{\mathbf{s} \in \mathcal{S}}\; F(\mathbf{s}) = -\mathrm{LCOE}(\mathbf{s}) - \sum_{c \in \mathcal{C}} \gamma_c \cdot \Phi\left(v_c(\mathbf{s}),\; \bar{v}_c\right) + 1 \times \delta_{v_c(\mathbf{s}) \leq \bar{v}_c}
\end{equation}

при ограничениях Г1--Г7,

где допустимое множество $\mathcal{S}$ --- совокупность всех четвёрок $(T, \sigma_0, \sigma_1, \sigma_2)$, удовлетворяющих ограничениям Г1--Г7. Размер $|\mathcal{S}|$ экспоненциально зависит от $N$, $M$, $N_0$ и на практике достигает $\sim 10^{100}$ и выше, что исключает точный перебор.

\textit{Замечание об асимметрии переменных.} Функция $\sigma_0$ не является инъекцией --- один вид свежей ТВС может стоять в нескольких позициях, поэтому ограничение Г2 мягче, чем Г5--Г6. Функции $\sigma_1$ и $\sigma_2$ --- строго инъективны, поскольку речь идёт о конкретных физических объектах, которые нельзя поставить в два места сразу. Это содержательное различие между «выбором вида» и «перестановкой конкретных объектов» является ключевой структурной особенностью данной задачи.


\end{document}
